% applications/use_cases-DE.tex
\chapter{Anwendungsfälle}
\label{chap:applications-use-cases}

Die Anwendungsfälle dieses Kapitels sind als Anwendungsfolgen genehmigter
Begriffe strukturiert.

\section{Struktur eines Anwendungsfalls}

Jeder Anwendungsfall folgt der Abfolge
\[
\text{Begriffskontext}
\rightarrow
\text{Ingenieuraufgabe}
\rightarrow
\substack{\text{Laufzeit-/Realisierungs-}\\\text{dienste}}
\rightarrow
\substack{\text{Ergebnis der}\\\text{Entscheidungsunterstützung}}.
\]

\section{Repräsentative Anwendungsfälle}

Repräsentative Anwendungsfälle umfassen die Rekonstruktion aus
heterogenen Eingaben, Glättung und Übergangsverfeinerung, gekoppelte
Krümmungs-Überhöhungs-Auswertung, Vergleiche realisierter Zustände sowie
die Interpretation von Mehrgleiskontexten.

Alle Anwendungsfälle verwenden kanonische Zustands-, Laufzeit- und
Realitätsbegriffe, anstatt sie neu zu definieren.

\section{Konsequenzen für Arbeitsabläufe}

Anwendungsabläufe verbinden Dateninterpretation, Laufzeitauswertung,
realisierungsbewussten Vergleich und Interoperabilitätsausgabe. Sie
unterstützen Ingenieurentscheidungen wie die Priorisierung der
Instandhaltung, die Abstimmung betrieblicher Bereiche und den Vergleich
von Entwurfsalternativen.

\section{Verbindung zu AIM}
\begin{summary}
Die Anwendungsfälle zeigen, dass AIM einen wiederverwendbaren
Ingenieurkern für praktische Arbeitsabläufe bereitstellt.

Anwendungsverhalten wird durch kanonische Operatoren und
Dienstekompositionen abgeleitet, nicht durch ad hoc entwickelte,
anwendungsspezifische Zustandsmodelle.
\end{summary}
